\section{The model as formalised}\label{sec:model}

Aggregate shocks $z$ take finitely many values and follow a Markov chain $\pi(z,z')$. Each household
has an idiosyncratic state $e$ in a finite set, moving with probabilities $p(z,z',e,e')$ that may
depend on today's and tomorrow's aggregate shock, and labour efficiency $\epsilon(e)$. Households
hold positions $a-b$ with $a\in[0,\abar]$, where $b\ge0$ is a borrowing limit, discount at
$\beta<1$, and receive the gross return $R(z,K,L)$ on their position and the wage $w(z,K,L)$ per
efficiency unit, where $K$ is mean (shifted) assets and $L$ aggregate labour. Cash on hand is
$R(a-b)+w\epsilon(e)$; saving $a'$ is feasible when
$0\le a'\le\min(\abar,\ R(a-b)+w\epsilon(e)-\underline c)$, with $\underline c\ge 0$ a consumption
floor that labour income net of debt service at the limit always affords. Krusell and Smith's
economy has $b=0$; results that need it say so (the monotone route of
Section~\ref{sec:recursive} and the capital lower bound below, since a higher return hurts
debtors). Period utility is continuous, increasing and strictly concave on
$[\underline c,\infty)$.

Prices are abstract bounded continuous functions; the Cobb--Douglas instance is
$R=1-\delta+\alpha A_zK^{\alpha-1}L^{1-\alpha}$, $w=(1-\alpha)A_zK^\alpha L^{-\alpha}$ with capital
and labour clamped to boxes, so that prices are bounded. Utility unbounded below at zero (logarithmic,
CRRA with $\gamma\ge1$) enters through a positive income floor and a consumption floor proved never
to bind (\lean{KSInada}). A household's problem along a path of prices is a stochastic dynamic
program on the tree of shock histories; its value function is the fixed point of a Blackwell
contraction, and strict concavity gives a unique optimal policy, jointly continuous in assets and the
price path.

\begin{definition}
A \emph{sequential equilibrium} from shock $z_0$ and distribution $\mu_0$ is a path $K$ of capital on
the tree of shock histories such that, at every node, capital equals mean assets of the distribution
generated by optimal saving along the price path of $K$ (with labour the mean efficiency of the
idiosyncratic marginal).
\end{definition}

\noindent
Two features of the formalisation should be kept in mind. The asset cap $\abar$ is a constraint;
for small capital shares it is proved irrelevant (Theorem~\ref{thm:cap}). And the clamping boxes of the Cobb--Douglas instance are part of the price
specification; where it matters we prove that the equilibrium stays inside them, so that prices are
the true marginal products.
